\documentclass[11pt]{article}
\usepackage[T1]{fontenc}
\usepackage[utf8]{inputenc}
\usepackage{amsmath,amssymb,amsthm,mathtools}
\usepackage{geometry,graphicx,booktabs,longtable,tabularx,array}
\usepackage{enumitem,needspace,xcolor,tikz,tcolorbox}
\usetikzlibrary{arrows.meta,positioning,calc,fit,backgrounds,decorations.pathreplacing,matrix,patterns}
\tcbuselibrary{breakable}
\usepackage{caption,hyperref,bookmark,cleveref,microtype}
\usepackage{komorebi-publication}
\renewcommand{\KomorebiTitle}{Komorebi Space}
\renewcommand{\KomorebiSubtitle}{Coherent realizations, obstructions, and descent}
\renewcommand{\KomorebiSeries}{Foundations of the Realization Calculus}
\renewcommand{\KomorebiKind}{Monograph}
\renewcommand{\KomorebiEdition}{Expanded monograph edition 1.5.0 / 12 October 2026}
\renewcommand{\KomorebiShortTitle}{Komorebi Space}
\renewcommand{\KomorebiProfile}{monograph}
\KomorebiSetup
\setcounter{tocdepth}{1}
\setcounter{secnumdepth}{2}
\setlist{nosep,leftmargin=1.7em,topsep=4pt}
\hypersetup{pdfkeywords={Komorebi Space,realization calculus,descent,obstruction,Lean,holonomy,closed range}}
\newtheorem{theorem}{Theorem}[section]
\newtheorem{proposition}[theorem]{Proposition}
\newtheorem{lemma}[theorem]{Lemma}
\newtheorem{corollary}[theorem]{Corollary}
\theoremstyle{definition}
\newtheorem{definition}[theorem]{Definition}
\newtheorem{example}[theorem]{Example}
\theoremstyle{remark}
\newtheorem{remark}[theorem]{Remark}
\numberwithin{equation}{section}
\newcommand{\R}{\mathbb R}
\newcommand{\Z}{\mathbb Z}
\newcommand{\Q}{\mathbb Q}
\newcommand{\N}{\mathbb N}
\newcommand{\K}{\mathbb K}
\newcommand{\V}{\mathcal V}
\newcommand{\W}{\mathcal W}
\newcommand{\id}{\operatorname{id}}
\newcommand{\im}{\operatorname{im}}
\newcommand{\ran}{\operatorname{ran}}
\newcommand{\Hom}{\operatorname{Hom}}
\newcommand{\Ob}{\operatorname{Ob}}
\newcommand{\Ext}{\operatorname{Ext}}
\newcommand{\Match}{\operatorname{Match}}
\newcommand{\Exact}{\operatorname{Exact}}
\newcommand{\dist}{\operatorname{dist}}
\newcommand{\Aut}{\operatorname{Aut}}
\newcommand{\Stab}{\operatorname{Stab}}
\newcommand{\Set}{\mathbf{Set}}
\newcommand{\Cat}{\mathbf{Cat}}
\newcommand{\norm}[1]{\left\lVert#1\right\rVert}
\newcommand{\abs}[1]{\left\lvert#1\right\rvert}
\newcommand{\source}[1]{\par\phantomsection\label{src:#1}{\footnotesize\sffamily\color{KomorebiMuted}Source concordance: \hyperlink{claim:#1}{#1}.}\par\smallskip}
\newcommand{\code}[1]{\texttt{\detokenize{#1}}}
\newcommand{\chap}[1]{\par\Needspace{10\baselineskip}\section{#1}}
\newcommand{\partnote}[2]{\clearpage\part{#1}\noindent\begin{minipage}{.92\linewidth}\large #2\end{minipage}\par\vspace{12pt}}
\newtcolorbox{insight}{breakable,colback=KomorebiSoft!55,colframe=KomorebiPine!50,boxrule=.45pt,arc=1pt,left=9pt,right=9pt,top=7pt,bottom=7pt}
\newtcolorbox{scopebox}{breakable,colback=KomorebiWarm,colframe=KomorebiOchre!55,boxrule=.4pt,arc=1pt,left=9pt,right=9pt,top=7pt,bottom=7pt}
\usepackage{float}
\newcommand{\fixedplate}[3]{\begin{figure}[H]\centering\resizebox{.96\linewidth}{!}{\input{figures/#1}}\caption{#2}\label{#3}\end{figure}}
\newcommand{\plate}[3]{\begin{figure}[!htbp]\centering\resizebox{.96\linewidth}{!}{\input{figures/#1}}\caption{#2}\label{#3}\end{figure}}
\tikzset{box/.style={draw=KomorebiPine!70,rounded corners=2pt,fill=KomorebiSoft!50,align=center,inner sep=8pt,font=\small},arr/.style={-{Latex[length=2.2mm]},draw=KomorebiPine,thick},thinarr/.style={-{Latex[length=1.8mm]},draw=KomorebiMuted},lab/.style={font=\small,align=center},dot/.style={circle,fill=KomorebiPine,inner sep=2.3pt}}
\renewcommand{\arraystretch}{1.18}
\raggedbottom
\widowpenalty=7000\clubpenalty=7000
\setlength{\emergencystretch}{3em}

% Compact only the contents entries, without reducing body type size.
\usepackage{etoolbox}
\makeatletter
\patchcmd{\l@section}{1.0em}{0.6em}{}{}
\makeatother

% Generated by checks/edition_stats.py from the current executed reports.
% Generated from current executed reports. Do not hand edit.
\newcommand{\EditionClaims}{80}
\newcommand{\EditionReferences}{211}
\newcommand{\EditionWrittenOnly}{0}
\newcommand{\EditionOriginalChecks}{14,224}
\newcommand{\EditionRecoveredChecks}{2,767}
\newcommand{\EditionRevisionChecks}{17,213}
\newcommand{\EditionExpansionChecks}{96,345}
\newcommand{\EditionDeepeningChecks}{13,817}
\newcommand{\EditionTotalChecks}{178,754}
\newcommand{\EditionRegularizationChecks}{34,388}
\newcommand{\EditionWrittenExtensions}{12}
\newcommand{\EditionFormalModules}{138}
\newcommand{\EditionAuditedDeclarations}{4,745}
\newcommand{\EditionNewTheorems}{54}
\newcommand{\EditionFormalTheorems}{1,507}
\newcommand{\EditionNewConstructions}{17}
\newcommand{\EditionReplays}{1}
\newcommand{\EditionNegativeControls}{7}


% Keep readable half-page figures with their surrounding exposition.
\renewcommand{\topfraction}{.88}
\renewcommand{\bottomfraction}{.80}
\renewcommand{\textfraction}{.10}
\renewcommand{\floatpagefraction}{.78}
\setcounter{topnumber}{2}
\setcounter{bottomnumber}{2}
\setcounter{totalnumber}{3}

% WR identifiers are proved-in-print extensions, not inherited Lean claims.
\newcommand{\written}[1]{\par\smallskip{\footnotesize\sffamily\color{KomorebiMuted}Written extension \hyperref[app:written]{#1}. No new Lean formalization is claimed.}\par\smallskip}

\begin{document}
\KomorebiCover
\pagenumbering{roman}
\section*{About this edition}
This monograph develops Komorebi Space as a theory of coherent realizations relative to a specified system of mathematical views. Its central constructions are ordinary limits, relative extension spaces, obstruction quotients, cocycle groupoids, an explicit stack on a composition site, and quantitative repair operators. The point of the synthesis is to keep these constructions connected without erasing the differences between their hypotheses.

The exposition follows the mathematics rather than the order in which the source library grew. Definitions precede existence questions. Existence is separated from uniqueness. Algebraic feasibility is separated from numerical stability. Changing a representation is separated from changing the data that the representation is required to preserve.

Publication version \textbf{1.5.0} develops the next step after the escaping Hilbert model. Four added chapters prove a bounded-correction existence criterion, strong convergence of minimum energy finite repairs, sharp fixed-budget certificates, and protected quadratic regularization. The diagonal model then gives an explicit asymptotic optimum and a contrasting attainable noisy problem. These twelve additions have complete printed proofs and distinct WR identifiers. They are not represented as new Lean declarations. The existing eighty source-linked KS entries retain their accepted formal baseline.

The canonical Lean baseline is version 1.2.0. It contains \EditionFormalModules\ production modules and one root. Its accepted clean replay and imported axiom audit cover \EditionAuditedDeclarations\ project declarations, including generated declarations. This edition retains those records and checks that every formal file is unchanged. It does \emph{not} claim a new Lean compilation. Fresh checks in this edition cover source correspondence, exact finite models, figure generation, and the publication build. Historical archives and compiler records retain their original dates and scopes.

Lean's compiler, kernel, and pinned external dependency cache remain trusted. This is not an independent rebuild of mathlib or a proof that a Python program implements every formal construction. Source correspondence identifies the actual declarations and hypotheses. It does not turn explanatory prose into a kernel-checked language.

All mathematical figures are reproducible vector diagrams or Python plots of explicit formulas. Axes, support patterns, and worked values are checked against their stated mathematical models. The source package contains the manuscript, figure sources, exact checks, evidence archives, and the progress ledger. The separate visual guide offers a shorter route through the examples. The Lean guide explains the correspondence and its verification boundaries.

\begin{scopebox}
\textbf{Reading scope.} The word \emph{space} does not assert one universal topological object containing all of mathematics. The strict construction is a set of coherent families. Symmetry is retained in a groupoid, and effective descent is proved for a specified site. Analytic estimates apply only after a normed or Hilbert model has been declared. These are related layers, not interchangeable definitions.
\end{scopebox}

\clearpage
\section*{Orientation}
A description at one scale often omits information visible at another. A formula in one representation may fail to commute with a change of representation. Local pieces may each be valid and still fail to assemble. Even when assembly is possible, there may be many completions, and small discrepancies may or may not imply proximity to an exact answer.

Komorebi Space treats these questions in a common order. First specify the views and the maps between them. Then define the coherent realizations. Next identify the information that must remain fixed. Finally determine the obstruction to extension, the freedom in a successful extension, and the conditions under which repair is stable.

The resulting program has three central distinctions:
\[
\begin{gathered}
\text{a collection of appearances}\quad\ne\quad\text{a coherent realization},\\
\text{existence of an exact answer}\quad\ne\quad\text{a stable repair estimate},\\
\text{one isomorphism class}\quad\ne\quad\text{absence of symmetries}.
\end{gathered}
\]
The monograph makes these distinctions into explicit constructions and theorems.

\subsection*{The route through the book}
Part I constructs coherent sections and follows their obstructions as more data are prescribed. Part II develops resolution towers, finite observations, split repair coordinates, and multiscale reconstruction. Part III retains symmetries, proves effective descent, and distinguishes coherent phases from removable ones. Part IV develops analytic stability and an independent geometry of transformation costs. Part V follows approximate answers to their limits, distinguishes a coherent point from its intended Hilbert realization, and then derives the bounded-budget and regularization criteria that diagnose the difference. Part VI works through two finite loop models and returns to mathematical semantics. Appendix~\ref{app:workshop} contains twelve solved problems. Appendix~\ref{app:hypotheses} gives a hypothesis atlas, and Appendix~\ref{app:written} records the exact evidence scope of the new written extensions.

A short technical route is Sections~\ref{sec:views}, \ref{sec:obstruction}, \ref{sec:stack}, \ref{sec:closed}, and~\ref{sec:motion}. A reader interested in the role of packets and residuals should include Sections~\ref{sec:split} and~\ref{sec:cascade}. The graph models in Section~\ref{sec:graph} explain how a categorical realization problem acquires an actual numerical condition number. Sections~\ref{sec:observation} and~\ref{sec:repair-transport} give the representation-theoretic links. Sections~\ref{sec:signal-loop}, \ref{sec:twisted-loop}, and~\ref{sec:escape} can also be read first as concrete motivations. The last of these connects finite realization, closure, and failure of protected infinite existence in one calculation. Sections~\ref{sec:bounded-realization} through~\ref{sec:diagonal-regularization} give its positive converse, fixed-budget diagnostic, and regularized continuation.

\subsection*{What is being contributed}
The limit universal property, gauge action groupoids, standard stack definitions, orthogonal projection, and bounded inverse principles belong to established mathematics. They are used here as such. The contribution of this presentation is their explicit integration with the project's actual view systems, protected correction spaces, resolution differentials, graph transports, and checked examples. Historical novelty of every component is not inferred from formal verification. Conversely, the fact that the ingredients are classical does not remove the need to prove the comparison maps and to retain the right hypotheses.

\clearpage\begingroup\setlength{\parskip}{0pt}\tableofcontents\endgroup
\clearpage\pagenumbering{arabic}
\input{chapters/01_views}
\input{chapters/02_obstructions}
\input{chapters/02a_prescription}
\input{chapters/03_towers}
\input{chapters/04_symmetry}
\input{chapters/04a_boundary_lab}
\input{chapters/04b_phase}
\addtocontents{toc}{\protect\newpage}
\input{chapters/05_geometry}
\input{chapters/05a_certificates}
\input{chapters/05b_transport_cost}
\input{chapters/05c_iteration}
\input{chapters/05d_floors}
\input{chapters/05e_escape}
\clearpage
\chap{A bounded repair budget can recover existence}\label{sec:bounded-realization}
The preceding example does not merely lack a convergent iteration. It has no bounded family of increasingly accurate protected corrections. This suggests a converse question. If approximate corrections stay in one bounded region, does an exact answer exist?

For bounded linear equations in a Hilbert space, the answer is yes. This is an existence statement, not a claim that the given sequence converges in norm. A stronger conclusion becomes available when the candidates are the least norm solutions of nested finite constraint systems.

\begin{scopebox}
This chapter and the next three are written extensions, indexed WR01--WR12. Their proofs are supplied here. They are not additional Lean declarations in the retained formal baseline. Their exact finite models have separate executable checks. The Hilbert and regularization background is classical \cite{clason-inverse}. No novelty is claimed for weak compactness or quadratic regularization.
\end{scopebox}

\subsection{The correction space is the protected kernel}
Let $E,Y,Z$ be real Hilbert spaces, let $d:E\to Y$ and $p:E\to Z$ be bounded linear maps, and fix a raw candidate $f\in E$. Write
\[
 H=\ker p,\qquad T=d|_H:H\to Y,\qquad r=df.
\]
The space $H$ is a closed Hilbert subspace with the norm inherited from $E$. An allowed correction is $a\in H$, and its repaired candidate is $f-a$. Thus
\[
 p(f-a)=pf,\qquad d(f-a)=r-Ta.
\]
The protected exact problem is precisely $Ta=r$. Restricting the domain before optimizing is important. A penalty on $\norm{pa}$ would only discourage a change in the protected data, not prohibit it.

The adjoint below always means $T^*:Y\to H$, taken with the inherited Hilbert metric on $H$. In ambient coordinates this is the projection of $d^*$ onto $H$. A coordinate basis of $\ker p$ need not be orthonormal. Its Gram matrix must then be included in any calculation of correction energy.

\subsection{The compactness needed for the existence argument}
\begin{lemma}[Bounded sequences have weakly convergent subsequences]\label{lem:weak-subsequence}
Every bounded sequence $(a_n)$ in a real Hilbert space has a subsequence converging weakly to some $a$. Moreover,
\[
 \norm a\le\liminf_k\norm{a_{n_k}}.
\]
\end{lemma}
\begin{proof}
All $a_n$ lie in their separable closed linear span $H_0$. In the infinite dimensional case, choose a countable orthonormal basis $(e_j)$ of $H_0$. Each scalar sequence $\langle a_n,e_j\rangle$ is bounded. Successive subsequences and a diagonal choice produce a subsequence on which every such coordinate converges, say to $c_j$. If $\norm{a_n}\le M$, Bessel's inequality on any finite set gives
\[
 \sum_{j=1}^{J}c_j^2
 =\lim_k\sum_{j=1}^{J}\langle a_{n_k},e_j\rangle^2\le M^2.
\]
Consequently $a=\sum_{j\ge1}c_je_j$ belongs to $H_0$ and has norm at most $M$. To test weak convergence against $h\in H_0$, approximate $h$ by a finite linear combination of the $e_j$. Coordinate convergence handles that finite part, while Cauchy--Schwarz bounds the remaining inner products uniformly by $2M$ times the approximation error. The component of any $h\in H_0^\perp$ pairs to zero with the entire sequence and its limit. The finite dimensional case uses the same argument with finitely many coordinates. The zero space is immediate.

Finally, when $a\ne0$, weak convergence tested against $a/\norm a$ gives $\norm a\le\liminf_k\norm{a_{n_k}}$. The inequality is automatic when $a=0$.
\end{proof}

\begin{theorem}[Bounded approximate feasibility]\label{WR01}
For a bounded linear map $T:H\to Y$ between real Hilbert spaces and $r\in Y$, the following are equivalent:
\[
 \begin{split}
 &r\in\ran T,\\
 &\text{there is a bounded sequence }(a_n)\subset H
       \text{ with }\norm{Ta_n-r}\longrightarrow0.
 \end{split}
\]
If all $\norm{a_n}\le M$, an exact solution exists with norm at most $M$.
\end{theorem}
\begin{proof}
An exact solution gives a constant sequence. In the other direction, take a weakly convergent subsequence $a_{n_k}\rightharpoonup a$. For every $y\in Y$,
\[
 \langle Ta,y\rangle
 =\langle a,T^*y\rangle
 =\lim_k\langle a_{n_k},T^*y\rangle
 =\lim_k\langle Ta_{n_k},y\rangle
 =\langle r,y\rangle.
\]
Thus $Ta=r$. The norm statement follows from the lemma.
\end{proof}
\written{WR01}

No closed range assumption appears. The theorem concerns a single datum and a bounded family of its corrections. Closed range is the stronger, uniform statement that every attainable datum admits a correction controlled by one residual gain.

An immediate contrapositive is useful. If $r\in\overline{\ran T}\setminus\ran T$, then every sequence with $Ta_n\to r$ satisfies $\norm{a_n}\to\infty$. Otherwise infinitely many terms would lie in one bounded ball, contradicting the theorem along that subsequence. This recovers the escape mechanism of Section~\ref{sec:escape} without using the diagonal formula.

\subsection{Existence does not make an arbitrary sequence converge strongly}
Take $H=Y=\ell^2$ and $Te_n=e_n/n$. The vectors $a_n=e_n$ have norm one and satisfy $Ta_n\to0$. An exact answer exists, namely zero. Yet $\norm{e_n-e_m}=\sqrt2$ for $n\ne m$, so no subsequence converges in norm.

The conclusion of WR01 is the existence of a weak limit along a subsequence that solves the equation. It is not summable movement, a strong convergence rate, or a selection rule for the original candidates. The minimum correction for the zero datum is zero, which this particular sequence never approaches in norm.

\subsection{Minimum energy finite repairs do converge strongly}
\fixedplate{31_nested_energy}{Minimum energy finite corrections behave differently for two data under the same diagonal operator. Bounded partial energy produces a strong Hilbert realization. Escaping partial energy accompanies a coherent product-space point outside the Hilbert space. The attainable energy lies below the explicit bound two.}{fig:nested-energy}

Suppose now that $A_1\supseteq A_2\supseteq\cdots$ are nonempty closed affine subsets of a real Hilbert correction space $H$. A typical $A_n$ consists of corrections satisfying the first $n$ bounded linear observation equations. The sets are affine because the equations are linear, and closed because the observations are continuous.

Let $a_n$ be the unique point of $A_n$ nearest to zero, and put $m_n=\norm{a_n}$. The Hilbert projection theorem supplies these points. It can also be proved directly: a sequence approaching the minimum norm is Cauchy by the parallelogram identity, its limit belongs to the closed convex set, and the same identity gives uniqueness.

\begin{theorem}[A uniform finite budget yields a strong realization]\label{WR02}
For the nested affine sets above, $m_n$ is nondecreasing, and for $m\ge n$,
\begin{equation}\label{eq:nested-energy}
 \norm{a_m-a_n}^2=m_m^2-m_n^2.
\end{equation}
Their intersection is nonempty if and only if $\sup_n m_n<\infty$. In that case $a_n$ converges in norm to the unique least norm point $a_*\in\bigcap_n A_n$, with
\[
 \norm{a_*-a_n}^2=\norm{a_*}^2-\norm{a_n}^2.
\]
If the intersection is empty, then $m_n\to\infty$.
\end{theorem}
\begin{proof}
Write $A_n=a_n+L_n$, where $L_n$ is a closed linear subspace. Minimality of $a_n$ on every line $a_n+t\ell$, $\ell\in L_n$, gives $\langle a_n,\ell\rangle=0$. Since $a_m-a_n\in L_n$ when $m\ge n$, the Pythagorean identity proves \eqref{eq:nested-energy}. Monotonicity also follows directly from $A_m\subseteq A_n$.

If $m_n$ is bounded, $m_n^2$ converges, and \eqref{eq:nested-energy} makes $(a_n)$ Cauchy. Completeness gives $a_*$. For each fixed $n$, the tail lies in the closed set $A_n$, so $a_*\in A_n$. Every $a\in\bigcap_n A_n$ has $\norm a\ge m_n$ for every $n$. Passing to the limit shows $\norm a\ge\norm{a_*}$. Uniqueness follows from strict convexity of the norm square. Letting $m\to\infty$ in \eqref{eq:nested-energy} proves the error formula.

Conversely, any point in the intersection bounds all $m_n$ from above. If the intersection is empty, the nondecreasing sequence $m_n$ must therefore be unbounded and hence tend to infinity.
\end{proof}
\written{WR02}

The finite answers need not extend one another coordinate by coordinate. Their shared feature is minimum energy in one common ambient Hilbert space. This is the extra geometry that forces their differences to be orthogonal to the earlier minimizers.

For the attainable diagonal datum $r_n=1/n^2$ under $Tu=(u_n/n)$, the first $N$ equations force $a_n=1/n$ for $n\le N$. Their minimum energy completion is zero beyond $N$, and
\[
 m_N^2=\sum_{n=1}^{N}\frac1{n^2}\le2.
\]
These corrections converge strongly to $(1/n)$. For the escaping datum $r_n=1/n$, the corresponding energy is $m_N^2=N$. Every finite equation is again solved, but the budget test now rules out an ambient solution.


\begin{insight}
Finite consistency is not enough by itself. For nested closed affine problems in one Hilbert space, a uniform bound on the minimum correction norms is exactly what upgrades finite feasibility to ambient feasibility. This is a datum-specific test, not a uniform inverse bound for the operator.
\end{insight}

\chap{The attainable residual at a fixed correction budget}\label{sec:budget-frontier}
The unrestricted floor $\dist(r,\ran T)$ forgets how large a correction must be. To retain that information, fix a nonnegative budget $M$ and define
\[
 \beta_{T,r}(M)=\inf_{\norm a\le M}\norm{r-Ta}.
\]
The norm is the actual correction norm in $H=\ker p$. It is not the norm of arbitrary coordinates used to represent that kernel. The function $\beta$ describes the smallest residual compatible with a declared amount of allowed change.

\subsection{At every finite budget the residual minimum exists}
\begin{theorem}[Attainment and shape of the budget profile]\label{WR03}
For a bounded linear map between real Hilbert spaces, the infimum defining $\beta(M)$ is attained for every $M\ge0$. The profile is nonincreasing and convex. For $0\le M_1\le M_2$,
\[
 0\le\beta(M_1)-\beta(M_2)\le\norm T(M_2-M_1).
\]
Writing $C=\overline{\ran T}$, $b=P_Cr$, and $q=\norm{r-b}$, one has
\[
 \lim_{M\to\infty}\beta(M)=q.
\]
The equality $\beta(M)=q$ at a finite budget is equivalent to the existence of $a$ with $Ta=b$ and $\norm a\le M$. In particular,
\[
 \beta(M)=0\quad\Longleftrightarrow\quad
 \exists a\in H:\ Ta=r,\ \norm a\le M.
\]
\end{theorem}
\begin{proof}
Choose a minimizing sequence in the ball of radius $M$. WR01's weak subsequence argument gives a weak limit in the same ball. The residuals converge weakly to $r-Ta$, and weak lower semicontinuity of the norm shows that this limit attains the minimum.

For convexity, mix minimizers at budgets $M_1,M_2$. Their convex combination has norm at most the same convex combination of the budgets, and the triangle inequality bounds its residual by the corresponding combination of residuals. Monotonicity follows from inclusion of the balls. For the Lipschitz estimate with $M_2>0$, scale a minimizer $a_2$ by $M_1/M_2$. The scaled correction is feasible at $M_1$ and differs from $a_2$ in norm by at most $M_2-M_1$. The case $M_2=0$ is immediate.

All $Ta$ lie in $C$. Hence
\begin{equation}\label{eq:range-decomposition}
 \norm{r-Ta}^2=q^2+\norm{b-Ta}^2.
\end{equation}
This bounds the profile below by $q$. Since $b$ is in the closure of the range, some finite correction achieves residual below any number strictly larger than $q$. That correction belongs to a sufficiently large budget ball, proving the limit. Finally, attainment and \eqref{eq:range-decomposition} give all the stated equality criteria.
\end{proof}
\written{WR03}

Although $\ran T$ may not be closed, the image $T(\{a:\norm a\le M\})$ of each closed budget ball is closed and convex. To see closedness, take a convergent sequence in that image, lift it to the bounded ball, and apply the weak limit argument. The nonclosed range is the increasing union of these closed images. Passing to an unlimited budget can therefore lose closedness even though every fixed-budget optimization is well posed.

The best residual vector is unique because it is the projection of $r$ onto a closed convex set. Its preimages can differ by vectors in $\ker T$. Among those preimages there is a unique least norm correction. These two uniqueness statements should not be confused with injectivity of $T$.

\subsection{A dual witness can rule out every correction in the budget}
\begin{theorem}[A sharp budget certificate]\label{WR04}
For every $M\ge0$,
\begin{equation}\label{eq:budget-dual}
 \beta_{T,r}(M)=
 \max_{\norm z\le1}\bigl(\langle r,z\rangle-M\norm{T^*z}\bigr).
\end{equation}
Thus any single $z$ with $\norm z\le1$ and
\[
 \langle r,z\rangle-M\norm{T^*z}>\varepsilon
\]
certifies that no correction of norm at most $M$ has residual at most $\varepsilon$.
\end{theorem}
\begin{proof}
For every feasible $a$ and every $\norm z\le1$,
\[
 \langle r,z\rangle-M\norm{T^*z}
 \le\langle r-Ta,z\rangle\le\norm{r-Ta}.
\]
This proves one inequality. Let $a_M$ minimize the residual, and set $v=r-Ta_M$. If $v=0$, the upper bound is zero and $z=0$ attains it. Otherwise let $z=v/\norm v$.

Minimality on the convex ball implies
\[
 \langle v,T(a-a_M)\rangle\le0
 \quad\text{for every }\norm a\le M.
\]
This follows by differentiating the residual square along the segment from $a_M$ to $a$ at its initial endpoint. Therefore $a_M$ maximizes $\langle T^*z,a\rangle$ on the ball. Its value is $M\norm{T^*z}$, including the cases $M=0$ and $T^*z=0$. Substitution gives
\[
 \langle r,z\rangle-M\norm{T^*z}
 =\langle r-Ta_M,z\rangle=\norm v=\beta(M).
\]
The maximum is thus attained by a specified witness, without an appeal to an interchange of infimum and supremum.
\end{proof}
\written{WR04}

In finite dimensions this is a checkable inequality, not merely an optimization prescription. A program may propose the correction or the witness. The displayed inequality can be verified independently of the proposal procedure. In a protected ambient calculation, use the restricted adjoint, not $d^*$ without projection onto $\ker p$.

\subsection{A two-dimensional calculation of the supporting witness}
Take
\[
 T=\begin{pmatrix}1&0\\0&1/2\end{pmatrix},\qquad r=(1,1),\qquad
 a=(1/2,2/5),\qquad v=r-Ta=(1/2,4/5).
\]
Then $T^*v=a$. At the budget $M=\norm a=\sqrt{41}/10$, the normalized vector $z=v/\norm v$ gives
\[
 \langle r,v\rangle-\norm a\norm{T^*v}
 =\frac{13}{10}-\frac{41}{100}=\frac{89}{100}=\norm v^2.
\]
Consequently $a$ is a globally optimal budgeted correction and $\beta(M)=\sqrt{89}/10$. It is not necessary to compare it with a dense numerical grid of competing vectors. The supporting witness proves optimality over the entire ball.

\fixedplate{32_budget_certificate}{The image of a Euclidean correction ball is an ellipse. The nearest reachable point has a residual normal to a supporting line. That same residual, normalized, is the dual witness in the exact budget formula. The plotted point and budget are the rational example in the text.}{fig:budget-certificate}

\subsection{How measurement and model uncertainty affect the budget}
\begin{theorem}[A robust budget profile]\label{WR05}
Suppose $\widetilde T:H\to Y$ and $\widetilde r\in Y$ satisfy
\[
 \norm{T-\widetilde T}\le\eta,\qquad
 \norm{r-\widetilde r}\le\delta,\qquad\eta,\delta\ge0.
\]
For every $M\ge0$,
\[
 \abs{\beta_{T,r}(M)-\beta_{\widetilde T,\widetilde r}(M)}
 \le\delta+\eta M.
\]
In particular, for $\norm z\le1$,
\[
 \beta_{T,r}(M)\ge
 \langle\widetilde r,z\rangle-M\norm{\widetilde T^*z}
       -\delta-\eta M.
\]
\end{theorem}
\begin{proof}
At any correction with $\norm a\le M$, the two residual vectors differ by
$(r-\widetilde r)-(T-\widetilde T)a$, of norm at most $\delta+\eta M$. The reverse triangle inequality bounds the difference between their norms by the same number. Take infima in each direction. Apply WR04 to the nominal system for the second assertion.
\end{proof}
\written{WR05}

This estimate identifies a practical boundary of increasingly large corrections. A model error $\eta$ is amplified by the correction budget even if the measured residual is tiny. The candidate-level inequality is
\[
 \norm{r-Ta}\le\norm{\widetilde r-\widetilde Ta}+\delta+\eta\norm a.
\]
The theorem compares maps on one fixed correction space. It does not cover an uncertain preservation map whose kernel itself changes.

\fixedplate{33_budget_uncertainty}{The nominal budget profile is surrounded by the rigorous uncertainty band $\delta+\eta M$. A large budget can improve the nominal fit while enlarging the uncertainty in that conclusion. The example uses $T=1$, $r=1$, $\delta=1/20$, and $\eta=1/10$.}{fig:budget-uncertainty}

\clearpage
\chap{Regularization selects a protected correction}\label{sec:regularized-repair}
A fixed budget asks how small the residual can become without exceeding an allowed correction size. Quadratic regularization reverses the emphasis. It asks for the best compromise between the residual and the size of the change. The resulting problem has a unique answer even when the exact protected equation has no solution.

That last sentence is both its usefulness and its limitation. The existence of a regularized answer is not a proof of exact realizability. The limiting behavior, rather than the mere existence of the minimizer, supplies the diagnosis.

\subsection{The normal equation is derived in the protected space}
For $\alpha>0$, define on $H=\ker p$
\[
 J_\alpha(a)=\norm{Ta-r}^2+\alpha\norm a^2.
\]
Both norms are the declared Hilbert norms. All minimizers below are corrections in $H$, so preservation is exact for every parameter, not just in the limit. This is ordinary quadratic Tikhonov regularization applied to the actual protected correction operator. The standard regularization background is discussed in \cite{clason-inverse}. The proof here does not require $T$ to be compact.

\begin{theorem}[A unique regularized correction]\label{WR06}
For every $\alpha>0$, there is a unique minimizer
\[
 a_\alpha=(T^*T+\alpha I)^{-1}T^*r.
\]
It satisfies $T^*(r-Ta_\alpha)=\alpha a_\alpha$. For every $a\in H$,
\begin{equation}\label{eq:regularization-energy}
 J_\alpha(a)-J_\alpha(a_\alpha)
 =\norm{T(a-a_\alpha)}^2+\alpha\norm{a-a_\alpha}^2.
\end{equation}
\end{theorem}
\begin{proof}
Let $B=T^*T+\alpha I$. Then $B$ is bounded and self-adjoint, and
\[
 \langle Ba,a\rangle=\norm{Ta}^2+\alpha\norm a^2
 \ge\alpha\norm a^2.
\]
Cauchy--Schwarz gives $\norm{Ba}\ge\alpha\norm a$. Thus $B$ is injective and has closed range: a convergent sequence $Ba_n$ forces $(a_n)$ to be Cauchy, and continuity identifies the limit. If $y$ is orthogonal to the range, self-adjointness gives $By=0$, hence $y=0$. The closed range is therefore dense and equals $H$. Its inverse exists and has norm at most $1/\alpha$.

Define $a_\alpha=B^{-1}T^*r$. Expanding $J_\alpha(a_\alpha+h)$, the two linear terms cancel by the normal equation, leaving \eqref{eq:regularization-energy}. The right side is nonnegative and can vanish only for $h=0$. This proves existence, global minimality, and uniqueness.
\end{proof}
\written{WR06}

A declared regularization parameter is part of the problem. A calculation that changes $\alpha$ is solving a different optimization problem, even if the protected data and forward operator stay fixed. On the other hand, the preservation condition itself does not change with $\alpha$.

\subsection{Coordinate bases must retain the correction metric}
In finite dimensions, let $B$ have linearly independent columns spanning $\ker p$. Write $a=Bc$. The coefficient vector $c$ is a representation of the correction, not necessarily an isometric one. Set
\[
 S=dB,\qquad G=B^{\mathsf T}B.
\]
The correct coordinate equation is
\begin{equation}\label{eq:gram-regularization}
 (S^{\mathsf T}S+\alpha G)c=S^{\mathsf T}r.
\end{equation}
The matrix $G$ is positive definite because the columns are independent. Replacing $G$ by the identity changes the meaning of correction size unless the chosen basis is orthonormal.

For example, let $d(a_1,a_2)=a_1$, $p(a_1,a_2)=a_1-a_2$, $r=1$, and $\alpha=1$. A basis of $\ker p$ is the column $B=(1,1)^{\mathsf T}$. Here $S=1$ and $G=2$, so the correction is $(1/3,1/3)$. Penalizing the coordinate with $G=1$ instead gives $(1/2,1/2)$. The latter is not the optimizer for the ambient Euclidean norm. The difference is a metric change, not a discrepancy between equivalent calculations.

The executable example in this edition constructs the nullspace basis over the rationals, retains its Gram matrix, solves \eqref{eq:gram-regularization}, and returns both the correction and an exact normal-equation residual. It also handles a zero-dimensional protected kernel without pretending that a nonzero correction exists.

\subsection{Parameter monotonicity and the fixed-budget problem}
\begin{proposition}[Regularized answers lie on the budget frontier]\label{WR07}
If $0<\alpha<\gamma$, then
\[
 \norm{a_\alpha}\ge\norm{a_\gamma},\qquad
 \norm{r-Ta_\alpha}\le\norm{r-Ta_\gamma}.
\]
At $M_\alpha=\norm{a_\alpha}$,
\[
 \beta(M_\alpha)=\norm{r-Ta_\alpha}.
\]
When $v_\alpha=r-Ta_\alpha\ne0$, its normalization is an optimizing witness in WR04 at that budget.
\end{proposition}
\begin{proof}
Write $A_\alpha=\norm{a_\alpha}^2$ and $R_\alpha=\norm{r-Ta_\alpha}^2$. Optimality at the two parameters gives
\[
 R_\alpha+\alpha A_\alpha\le R_\gamma+\alpha A_\gamma,
 \qquad
 R_\gamma+\gamma A_\gamma\le R_\alpha+\gamma A_\alpha.
\]
Adding proves $(\gamma-\alpha)(A_\alpha-A_\gamma)\ge0$, and substitution into the first inequality gives $R_\alpha\le R_\gamma$.

If a correction with norm at most $M_\alpha$ had a strictly smaller residual, its $J_\alpha$ value would also be smaller. Thus $a_\alpha$ attains the budget minimum. The normal equation gives
\[
 T^*v_\alpha=\alpha a_\alpha,\qquad
 \langle r,v_\alpha\rangle-\alpha\norm{a_\alpha}^2
 =\norm{v_\alpha}^2.
\]
Divide the second equation by $\norm{v_\alpha}$ to obtain equality in the dual formula.
\end{proof}
\written{WR07}

This links a penalized calculation to a hard-budget certificate. It does not say that every larger budget must produce a different answer. Once the least-squares floor is attained, excess budget can be irrelevant. If $T^*r=0$, every regularized correction is zero and the entire profile is constant.

\subsection{The limit separates projection from actual realizability}
Keep $C=\overline{\ran T}$, $b=P_Cr$, and $q=\norm{r-b}$. Since $r-b\in\ker T^*$, the regularized corrections depend on $b$, not on the component that is orthogonal to the closure of all attainable residuals.

\begin{theorem}[The regularization limit alternative]\label{WR08}
As $\alpha\downarrow0$,
\[
 Ta_\alpha\longrightarrow b,\qquad
 \norm{r-Ta_\alpha}\longrightarrow q,\qquad
 \alpha\norm{a_\alpha}^2\longrightarrow0.
\]
If $b\in\ran T$, let $a^\dagger$ be its unique minimum norm preimage. Then
\[
 a_\alpha\longrightarrow a^\dagger\quad\text{in norm}.
\]
If $b\notin\ran T$, then $\norm{a_\alpha}\to\infty$.
\end{theorem}
\begin{proof}
For any fixed $a\in H$, minimality gives
\[
 q^2\le\norm{r-Ta_\alpha}^2
 \le J_\alpha(a_\alpha)\le\norm{r-Ta}^2+\alpha\norm a^2.
\]
Choose $a$ with residual square arbitrarily close to $q^2$, then let $\alpha\downarrow0$. It follows that both the residual square and the minimum objective tend to $q^2$. Their difference is $\alpha\norm{a_\alpha}^2$, so that term tends to zero. The orthogonal identity \eqref{eq:range-decomposition} then yields $Ta_\alpha\to b$.

If $b\in\ran T$, its preimage set is a nonempty closed affine space. Projection of zero gives $a^\dagger$. Comparison with $a^\dagger$ yields
\[
 \norm{Ta_\alpha-b}^2+\alpha\norm{a_\alpha}^2
 \le\alpha\norm{a^\dagger}^2,
\]
so $\norm{a_\alpha}\le\norm{a^\dagger}$. Along any sequence of parameters tending to zero, a weakly convergent subsequence has an exact preimage of $b$ as its limit. The norm bound and weak lower semicontinuity make that limit the minimum norm preimage $a^\dagger$. Along that subsequence, the norms converge to $\norm{a^\dagger}$, and weak convergence together with convergence of norms gives strong convergence. If the whole parameter family failed to converge strongly, a sequence staying a fixed distance away would have just such a strongly convergent subsequence, a contradiction.

If $b\notin\ran T$ and the correction norms failed to tend to infinity, some parameter sequence tending to zero would have bounded corrections. Their images tend to $b$, contradicting WR01.
\end{proof}
\written{WR08}

The vanishing penalty $\alpha\norm{a_\alpha}^2$ is not a boundedness certificate for $a_\alpha$. Its coefficient is tending to zero. In the escaping case, the penalty and residual can both tend to zero while the actual corrections diverge.

\begin{center}
\begin{tabularx}{\linewidth}{@{}p{.17\linewidth}p{.26\linewidth}X@{}}
\toprule Residual floor & Is $b$ attained? & Limit of the protected corrections\\\midrule
$q=0$ & Yes & Strong convergence to the nearest exact protected correction.\\[3pt]
$q=0$ & No & Residual vanishes, correction norms diverge, no exact protected answer.\\[3pt]
$q>0$ & Yes & Strong convergence to the least norm least-squares correction, not an exact answer.\\[3pt]
$q>0$ & No & A positive residual floor and diverging correction norms occur together.\\\bottomrule
\end{tabularx}
\end{center}

The last case is not hypothetical. Let $Tu=(u_n/n)$ on $\ell^2$, define $Su=(Tu,0)$ into $\ell^2\oplus\R$, and take $r=((1/n),1)$. Then $b=((1/n),0)$ is in the closure but not the range, while $q=1$. The same escaping corrections as before coexist with an irreducible orthogonal residual. Positive floor and unbounded correction cost are independent features.

\subsection{Noise is controlled at a fixed parameter}
\fixedplate{34_noise_filter}{For one singular value $\sigma$, the response to residual-data noise is $\sigma/(\sigma^2+\alpha)$. Its maximum occurs at $\sigma=\sqrt\alpha$ and equals $1/(2\sqrt\alpha)$. The plot evaluates this formula and the universal horizontal bounds at three explicit positive parameters.}{fig:noise-filter}

\begin{theorem}[Data sensitivity and a convergent parameter choice]\label{WR09}
For two residual data $r,\widetilde r$ with $\norm{r-\widetilde r}\le\delta$ and the same $T,\alpha$,
\[
 \norm{a_\alpha(r)-a_\alpha(\widetilde r)}
 \le\frac{\delta}{2\sqrt\alpha},\qquad
 \norm{T(a_\alpha(r)-a_\alpha(\widetilde r))}\le\delta.
\]
The coefficient $1/(2\sqrt\alpha)$ is sharp over bounded linear Hilbert operators. If $b=P_Cr\in\ran T$, $\alpha(\delta)\to0$, and $\delta/\sqrt{\alpha(\delta)}\to0$, then the noisy corrections converge strongly to $a^\dagger$ as $\delta\to0$.
\end{theorem}
\begin{proof}
Let $h=a_\alpha(\widetilde r)-a_\alpha(r)$ and $e=\widetilde r-r$. Subtracting the normal equations and pairing with $h$ gives
\[
 \norm{Th}^2+\alpha\norm h^2=\langle e,Th\rangle
 \le\delta\norm{Th}.
\]
Put $t=\norm{Th}$. Then $t\le\delta$ and
\[
 \alpha\norm h^2\le\delta t-t^2
 =\frac{\delta^2}{4}-(t-\delta/2)^2\le\frac{\delta^2}{4}.
\]
This proves the two estimates. For the scalar operator $T=\sqrt\alpha\,I$, the solution map is multiplication by $1/(2\sqrt\alpha)$, which proves sharpness. Finally, split the noisy error into the data perturbation at the chosen parameter and the exact-data regularization error. The first tends to zero by the displayed estimate and the second by WR08.
\end{proof}
\written{WR09}

The theorem does not manufacture an exact solution when $r$ is outside the range. For the positive-floor attainable case it converges to a least-squares correction. For the unattainable projected datum it cannot produce a bounded exact-data limit. The range hypothesis must therefore precede the parameter-choice conclusion.


\begin{insight}
Regularization gives a unique, exactly protected answer at every positive parameter. Whether those answers converge to an exact realization, converge only to a best approximation, or escape is determined by the projected datum's actual range membership. A smooth regularized path does not erase that distinction.
\end{insight}

\chap{A complete regularization diagnostic for the escaping model}\label{sec:diagonal-regularization}
The abstract alternatives become explicit for the diagonal model. They also answer a natural objection to the hard prefix corrections: perhaps a smoother correction could fit increasingly well without escaping. It can fit more efficiently, but it cannot avoid the obstruction.

\subsection{The entire regularized family is explicit}
Use $Tu=(u_n/n)$ and $r=w=(1/n)$ as in Section~\ref{sec:escape}. Choose $\alpha=m^{-2}$ for an integer $m\ge1$. The normal equation in each coordinate gives
\[
 a^{(m)}_n=\frac{m^2}{m^2+n^2},\qquad
 v^{(m)}_n=r_n-(Ta^{(m)})_n=\frac{n}{m^2+n^2}.
\]
The repaired ambient candidate is $(-a^{(m)},1)$. Its protected coordinate is one for every $m$.

\begin{proposition}[Energy, residual, and certified tails]\label{WR10}
The displayed $a^{(m)}$ belongs to $\ell^2$ and is the unique regularized correction at $\alpha=m^{-2}$. Set
\[
 A_m=\norm{a^{(m)}}^2=\sum_{n\ge1}\frac{m^4}{(m^2+n^2)^2},\qquad
 B_m=\norm{v^{(m)}}^2=\sum_{n\ge1}\frac{n^2}{(m^2+n^2)^2}.
\]
Then
\begin{equation}\label{eq:soft-bounds}
 \frac m4\le A_m\le\frac{4m}{3},\qquad
 \frac1{4(m+1)}\le B_m\le\frac2m.
\end{equation}
For the sums $A_{m,K},B_{m,K}$ through any integer $K\ge m$,
\begin{equation}\label{eq:soft-tails}
 0\le A_m-A_{m,K}\le\frac{m^4}{3K^3},\qquad
 0\le B_m-B_{m,K}\le\frac1K.
\end{equation}
\end{proposition}
\begin{proof}
For $n>m$, the correction square is at most $m^4/n^4$, whose sum converges. The residual belongs to $\ell^2$ because $T$ is bounded and $r\in\ell^2$. The coordinate normal equation is
\[
 \frac1n\frac{n}{m^2+n^2}=m^{-2}\frac{m^2}{m^2+n^2}.
\]
It holds in $\ell^2$, so WR06 proves global optimality.

For $n\le m$, the correction coordinate is at least $1/2$, yielding $A_m\ge m/4$. Bound its first $m$ squares by one and use
$\sum_{n>m}n^{-4}\le\int_m^\infty x^{-4}\,dx=1/(3m^3)$ for the tail. This gives $A_m\le m+m/3$.

The residual square on $n\le m$ is at most $n^2/m^4$, whose sum is at most $1/m$. The remaining sum is at most $\sum_{n>m}n^{-2}\le1/m$. For its lower bound, $n>m$ implies $(m^2+n^2)^2\le4n^4$, so
\[
 B_m\ge\frac14\sum_{n>m}\frac1{n^2}\ge\frac1{4(m+1)}.
\]
The same comparisons beyond $K$ give \eqref{eq:soft-tails}. All tail bounds concern the actual infinite sums.
\end{proof}
\written{WR10}

Thus the correction norm grows on the order of $\sqrt m$, the residual norm decays on the order of $m^{-1/2}$, and the penalty $\alpha A_m=A_m/m^2$ tends to zero. The smooth correction avoids the hard discontinuity at a prefix boundary, but it still approaches the forbidden constant sequence coordinate by coordinate.

\fixedplate{35_regularized_escape}{At each fixed coordinate, the soft correction approaches one as $m$ tends to infinity. Its transition moves outward as $m$ increases. The chart compares three exact coordinate formulas $m^2/(m^2+n^2)$, not finite approximations to an unspecified optimization procedure. Every displayed sequence has finite energy, but the family's energies are unbounded.}{fig:regularized-escape}

\subsection{The sharp asymptotic correction cost}
The preceding bounds establish divergence and decay. The leading constant can also be computed, which permits a meaningful comparison with hard prefix repair.

\begin{theorem}[A sharp budget asymptotic for the diagonal datum]\label{WR11}
As $m\to\infty$,
\[
 \frac{A_m}{m}\longrightarrow\frac\pi4,\qquad
 mB_m\longrightarrow\frac\pi4.
\]
For the exact budget profile of this infinite problem,
\begin{equation}\label{eq:sharp-budget}
 \lim_{M\to\infty}M\beta_{T,w}(M)=\frac\pi4.
\end{equation}
The hard prefix correction at budget $M=\sqrt N$ instead has
\[
 \sqrt N\,\norm{w-T\mathbf1_N}\longrightarrow1.
\]
\end{theorem}
\begin{proof}
The two normalized sums are Riemann sums on the positive half line:
\[
 \frac{A_m}{m}=\frac1m\sum_{n\ge1}\frac1{(1+(n/m)^2)^2},\qquad
 mB_m=\frac1m\sum_{n\ge1}\frac{(n/m)^2}{(1+(n/m)^2)^2}.
\]
On every bounded interval, ordinary Riemann convergence applies. Beyond $L>1$, the first function is bounded by $x^{-4}$ and the second by $x^{-2}$. Their sum tails are bounded uniformly for sufficiently large $m$ by constants times $L^{-3}$ and $L^{-1}$, respectively. The corresponding integral tails have the same bounds. Letting $m\to\infty$ first on the bounded interval and then $L\to\infty$ proves convergence of the full sums.

The substitution $x=\tan\theta$ gives
\[
 \int_0^\infty\frac{dx}{(1+x^2)^2}
 =\int_0^{\pi/2}\cos^2\theta\,d\theta=\frac\pi4,
\]
\[
 \int_0^\infty\frac{x^2\,dx}{(1+x^2)^2}
 =\int_0^{\pi/2}\sin^2\theta\,d\theta=\frac\pi4.
\]
Put $M_m=\sqrt{A_m}$. These budgets increase strictly to infinity, and $M_{m+1}/M_m\to1$ by the first asymptotic. WR07 gives $\beta(M_m)=\sqrt{B_m}$, so
$M_m\beta(M_m)=\sqrt{A_mB_m}\to\pi/4$.
For $M_m\le M\le M_{m+1}$, monotonicity of $\beta$ yields
\[
 M_m\beta(M_{m+1})\le M\beta(M)\le M_{m+1}\beta(M_m).
\]
Both bounds tend to $\pi/4$, proving \eqref{eq:sharp-budget}. Finally, the earlier telescoping tail enclosure gives
\[
 \frac{N}{N+1}\le N\sum_{n>N}\frac1{n^2}\le1,
\]
which proves the hard-prefix limit.
\end{proof}
\written{WR11}

At large matched correction budgets, the optimal soft repair therefore has about $\pi/4$ times the hard-prefix residual. This improves the efficiency of approximation, not the existence result. Both residuals tend to zero only as the budget tends to infinity. The limiting constant is a calculation for this declared model, not a universal regularization efficiency factor.

\fixedplate{36_budget_product}{The plotted intervals enclose the actual infinite quantities $\sqrt{A_mB_m}$, using finite sums plus the proved tails. The corresponding hard-prefix intervals enclose $\sqrt{N\sum_{n>N}n^{-2}}$ with $N=m$. These have different natural budget sequences. Their limits give the matched-budget asymptotic constants $\pi/4$ and one.}{fig:budget-product}

\subsection{The same operator has an attainable datum}
Change only the residual to $b_n=1/n^2$. Its minimum norm preimage is $a^\dagger_n=1/n\in\ell^2$. At $\alpha=m^{-2}$, the regularized correction is
\[
 c^{(m)}_n=\frac{m^2}{n(m^2+n^2)},\qquad
 a^\dagger_n-c^{(m)}_n=\frac{n}{m^2+n^2}.
\]
The bias square is exactly the previously computed $B_m$. This links the two cases without conflating their existence hypotheses: the escaping residual of one problem is the regularization bias of an attainable problem.

\begin{corollary}[An explicit noisy attainable repair]\label{WR12}
For the attainable datum $b=(1/n^2)$, suppose $\norm{\widetilde b-b}\le\delta$ with $0<\delta\le1$. Choose
\[
 m=\lceil\delta^{-2/3}\rceil,\qquad\alpha=m^{-2}.
\]
Then the noisy regularized correction satisfies
\[
 \norm{a_\alpha(\widetilde b)-a^\dagger}
 \le\sqrt{\frac2m}+\frac{\delta m}{2}
 \le(\sqrt2+1)\delta^{1/3}.
\]
\end{corollary}
\begin{proof}
The exact-data bias has norm $\sqrt{B_m}\le\sqrt{2/m}$ by WR10. The noise bound in WR09 is $\delta m/2$. The triangle inequality gives the first estimate. The ceiling choice obeys
$\delta^{-2/3}\le m\le\delta^{-2/3}+1\le2\delta^{-2/3}$, which gives the second.
\end{proof}
\written{WR12}

This is a deterministic guarantee for the specified datum and Hilbert noise bound. It is not a statistical confidence statement. The identical parameter rule in the escaping problem cannot produce convergence to an exact protected Hilbert solution, because no such solution exists.

\subsection{What the calculation permits a machine to certify}
The finite program can check a normal equation, exact preservation, and a quadratic optimality identity. It can return a dual witness ruling out a smaller residual at a given budget. In the infinite diagonal example it can return partial sums together with rational tail enclosures. Those are distinct certificates with distinct conclusions.

None of them alone establishes unrestricted semantic realizability. The bounded-budget theorem needs a genuinely uniform bound. The exactification theorem needs control of candidate movement. The regularization limit needs actual attainment of the projected datum. The repeated appearance of a small residual does not replace any of those hypotheses.

The useful next diagnostic is therefore not just ``Did the residual decrease?'' It is ``What happened to the correction norm, what quantity was preserved, and what existence theorem applies to this datum?''

\input{chapters/06_cases}
\chap{Mathematical meaning and the remaining boundaries}\label{sec:meaning}
\subsection{A relative theory can be complete without being absolute}
The strict space \(\Gamma(F)\) is canonical relative to its diagram: its universal property determines it up to the appropriate unique isomorphism. The additive obstruction is canonical relative to its admissible corrections. A marked realization problem is invariant under the specified comparison equivalences. A repair can be canonical relative to a protected observable, or orthogonal relative to a fixed inner product.

These are meaningful and different kinds of canonicity. They need not culminate in one object independent of every choice of views, coefficients, topology, semantics, and geometry. Some changes of view intentionally forget information. Others strengthen the question. Treating every such change as an equivalence would remove the very distinctions the calculus is designed to preserve.

The appropriate central statement is therefore not that every possible mathematical presentation lives in one fixed universal space. It is that a declared system of representations can be organized into coherent realizations, with explicit extension, comparison, and repair problems. A proposed bridge between two such systems has to show what structure it preserves and what information it loses.

\subsection{Internal realization and intended objects}
Suppose a separate category \(\mathcal S\) contains the intended semantic objects of an application, and a realization construction gives a comparison functor
\[
 \mathcal R:\mathcal S\longrightarrow\mathcal K,
\]
where \(\mathcal K\) is the relevant category or groupoid of coherent data. Defining \(\mathcal K\) does not establish that \(\mathcal R\) is faithful, full, or essentially surjective. These are new mathematical questions about the actual semantic comparison.

Faithfulness asks whether distinct semantic arrows can become indistinguishable. Fullness asks whether every compatible comparison comes from a semantic one. Essential surjectivity asks whether every intended coherent datum is realized by an actual semantic object. Depending on the application, one may seek only a restricted image or a classified obstruction to belonging to it.

This distinction prevents a formal globalization of encoded data from being mistaken for an unrestricted arithmetic or physical existence theorem. The abstract construction does not manufacture missing semantic objects.

\subsection{Every finite test can pass while the infinite problem fails}
\begin{example}[Finite prefixes without a global point]\label{KS54}
Let \(X_n=\{m\in\N:m\ge n\}\), with bonding maps given by inclusion. Every level is nonempty, and every finite prefix has a coherent point. The inverse limit is nevertheless empty.
\end{example}
\begin{proof}
For a prefix ending at \(N\), assign the integer \(N\) at every level. All inclusions preserve it. In a global point, however, the inclusion equations force every coordinate to be the same integer \(m\), while membership at level \(n\) requires \(m\ge n\) for every \(n\). Taking \(n=m+1\) is a contradiction.
\end{proof}
\source{KS54}
This counterexample is formalized in \code{FinitePrefixBoundary.lean} using the original resolution and partial-section types. It explains why inverse-limit existence needs more than separate finite solvability. Finiteness, compactness, surjectivity, or another coherence principle must be proved in the relevant application. The example does not contradict the strict gluing theorem: a compatible matching family of all prefixes was never supplied.

\subsection{A positive infinite-existence criterion}
The failed example isolates a missing compactness or lifting hypothesis. Finite fibers supply a concrete positive criterion without requiring every bonding map to be surjective.

\begin{proposition}[Finite fibers and persistent images]\label{KS58}
Let \((X_n,\rho_n)\) be a sequential inverse system of nonempty finite sets. Its inverse limit is nonempty. More precisely, if \(\rho_{n,m}:X_m\to X_n\) is composite restriction, then
\[
 \im(\operatorname{ev}_n)=Y_n,
 \qquad Y_n=\bigcap_{m\ge n}\im\rho_{n,m}.
\]
Thus a prescribed value at level \(n\) extends precisely when it survives every later image.
\end{proposition}
\begin{proof}
The images in the finite set \(X_n\) form a decreasing sequence of nonempty subsets, so they stabilize to a nonempty \(Y_n\). The maps \(\rho_n\) send \(Y_{n+1}\) into \(Y_n\). They are onto these subsets: given \(y\in Y_n\), choose a later stage \(m\) after both image sequences at levels \(n\) and \(n+1\) have stabilized. Lift \(y\) to \(X_m\) and restrict that lift to \(X_{n+1}\). It belongs to \(Y_{n+1}\) and maps to \(y\).

Choose successive preimages in these nonempty finite fibers. They give a coherent family. Starting at any specified \(y\in Y_n\), and restricting to earlier levels, proves the claimed evaluation image. The reverse inclusion follows by restricting a global family from each later stage.
\end{proof}
\source{KS58}
This is the standard finite-set inverse-limit principle, here with its sequential persistent-image proof~\cite{stacks-limits}. Its formalization in \code{FiniteRealizationExistence.lean} constructs an explicit equivalence with the pinned library's section type, then applies the finite cofiltered-system and eventual-range theorems. The printed sequential proof and this formal route establish the same stated criterion. Both use ordinary classical choice. Surjective bonding maps on arbitrary nonempty sets give a second sufficient condition by direct recursive lifting. Neither criterion says that arbitrary data in a non-surjective finite tower extend.

\subsection{A finite computation that exposes the preservation contract}
Consider a raw vector \(x=(2,0,5)\in\R^3\) with defect \(d(x)=x_1-x_2\). Unconstrained orthogonal repair gives \((1,1,5)\), with squared error two. If the first and third coordinates must be preserved, the unique exact repair is \((2,2,5)\), with squared error four.

Both outputs are exact. Only one satisfies the second preservation problem. Neither theorem contradicts the other because the admissible changes differ. An implementation that reports only zero residual would not distinguish the two tasks.

For a finite linear system over a field, the same question admits independently checkable positive and negative certificates. Let \(A\in\K^{m\times n}\) encode compatibility, \(P\in\K^{p\times n}\) the protected values, and \(x\in\K^n\) the raw candidate. Let the columns of \(N\) be a basis of \(\ker P\). The correction equation is
\[
 ANz=Ax,\qquad y=x-Nz.
\]
It is equivalent to the stacked system \(Ay=0\), \(Py=Px\).

\begin{proposition}[A finite certificate alternative]\label{KS57}
Exactly one of the following occurs:
\begin{enumerate}
\item There is \(y\) with \(Ay=0\) and \(Py=Px\). Every answer is \(y+H\), where \(H=\ker A\cap\ker P\).
\item There are row vectors \(\lambda\in\K^{1\times m}\) and \(\mu\in\K^{1\times p}\) with
\[
 \lambda A=\mu P,\qquad\lambda Ax=1.
\]
\end{enumerate}
The negative witness annihilates every admissible correction defect, since \(\lambda AN=\mu PN=0\). The ambiguity dimension in the positive case is
\[
 \dim H=n-\operatorname{rank}\begin{pmatrix}A\\P\end{pmatrix}.
\]
\end{proposition}
\begin{proof}
Row-reduce \(T=\begin{pmatrix}A\\P\end{pmatrix}\) with right side \(b=\begin{pmatrix}0\\Px\end{pmatrix}\), retaining the invertible row-operation matrix. A consistent reduced system supplies an answer, and differences of answers are exactly \(\ker T\). Rank-nullity gives the dimension.

Otherwise a row combination \(w=(u,v)\) satisfies \(wT=0\) and \(wb=\alpha\ne0\). Put \(\lambda=-u/\alpha\), \(\mu=v/\alpha\). Then \(\lambda A=\mu P\) and \(\lambda Ax=\mu Px=1\). Conversely an exact protected answer would imply
\[
 1=\lambda Ax=\mu Px=\mu Py=\lambda Ay=0,
\]
which is impossible in a field. Therefore the alternatives are exclusive and exhaustive.
\end{proof}
\source{KS57}
The formal proof gives these literal row-vector equations over an arbitrary field. It uses algebraic dual separation of the stacked image rather than verifying the row-reduction program. The affine answer equivalence and rank-nullity formula are proved separately. Thus a normalized obstruction is an actual checkable witness, while universal correctness of the supplied Python elimination implementation remains a different question.

\subsection{One equation, three preservation problems}
For the following three cases work over \(\Q\), with \(A=(1,-1,0)\) and \(x=(2,0,5)\). With \(P(x)=(x_1,x_3)\), the unique answer is \((2,2,5)\). With \(P(x)=x_3\), the entire answer set is \(\{(t,t,5):t\in\Q\}\). With \(P(x)=(x_1,x_2)\), no answer exists. In this last case
\[
 \lambda=(1/2),\qquad\mu=(1/2,-1/2)
\]
is a normalized negative certificate over \(\Q\) or \(\R\): \(\lambda A=\mu P\) and \(\lambda Ax=1\).

The package implements this finite interface over \(\Q\) in \code{examples/protected\_repair.py}. It returns a particular exact answer and an ambiguity basis, or a normalized dual witness. Its verifier checks the certificate equations directly. A Gaussian particular answer is \emph{not} advertised as the nearest answer of Theorem~\ref{KS56}. That is a separate geometric selection.

Exact field elimination does not solve the integer parity problem by changing its coefficient domain. Over \(\Q\), \(2m=1\) has the solution \(1/2\). Over \(\Z\), it has none. The rational interface rejects floating-point inputs rather than silently treating an approximate decimal as exact arithmetic. Its finite tests establish the tested implementation behavior, not an arbitrary-input Lean refinement theorem.

\subsection{Why this is relevant to machine reasoning}
The framework gives a machine explicit mathematical objects to manipulate: views, restrictions, partial realizations, typed defects, admissible corrections, and comparison witnesses. A successful step can carry both an output and a certificate of what remained fixed. A failed step can return an obstruction in a specified quotient rather than an unexplained failure flag.

The split coordinate theorem is especially direct. In its stated class, the same candidate is represented by protected information and residual data, and repair zeros only the residual coordinate. The descent theorem explains which local certificates can be assembled. The analytic results explain when a small measured defect has a quantitatively reliable meaning.

These observations motivate a machine oriented mathematical language. They do not establish an autonomous reasoning system, an empirical intelligence gain, or a general efficient algorithm. Classical choice, infinite products, quotient membership, and bounded inverse constructions can be mathematically valid without yielding executable finite procedures. Implementation claims need their own specifications and tests.

\subsection{What this monograph closes, and what it leaves open}
The book supplies a coherent account of strict realization spaces, additive and nonabelian extension problems, split tower repair, retained multiscale detail, a genuine composition-site stack, presentation transport, graph holonomy, closed-range stability, and moving exact spaces. The additional convergence and obstruction development separates small measurements from small movements, coherent coordinate families from their semantic image, and zero residual from zero transformation cost. The new bounded-budget development supplies a positive existence criterion and keeps range membership separate from regularized solvability. Its fixed-budget dual witness can certify a limit on every admissible correction, while its diagonal example computes the actual leading cost of approximation. Each layer is attached to the data that make it meaningful.

The following broader targets are not conclusions of this edition: arbitrary higher categorical coefficients and infinity-stack descent, unrestricted changes of view doctrine, norm reconstruction for every bounded infinite residual hierarchy, nonlinear or unknown-connection repair, perturbation across changing Hilbert geometries without new comparison data, and semantic essential surjectivity for unrestricted arithmetic or physical applications.

The integrated proof infrastructure for the stated source closure is supplied in this edition. The canonical project has \EditionFormalModules\ production modules, a single root, clean compiler records, declaration axiom reports, and rejection controls. All five previously identified written-only items KS54--KS58 have formal counterparts. That completed list is distinct from the newly developed WR01--WR12 results, whose full printed proofs have not been formalized in this edition. Further historical material is not silently included in this closure. Extensions of the theory require new hypotheses, sources, and verification records.

\begin{insight}
\textbf{The organizing principle.} Coherent realization, obstruction to extension, freedom of completion, and stability of repair belong to one mathematical conversation. Their unification is useful precisely when it does not identify them with one another.
\end{insight}

\appendix
\clearpage
\chap{A reader's workshop}\label{app:workshop}
These short problems isolate choices that recur in the main development. Each has a complete solution. They are reading exercises and direct applications of the cited theory, not additional entries in the formal theorem inventory.

\subsection*{1. A stable limit does not make every seed persistent}
At every level take \(X_n=\{0,1\}\), with each bonding map sending both elements to zero. Find the inverse limit and the image of evaluation at every level. Does nonemptiness of every finite fiber make the seed one extendable?

\textbf{Solution.} For a coherent family, \(x_n=\rho_n(x_{n+1})=0\) for every \(n\). The all-zero family is coherent and is therefore the only point. Every evaluation image is \(\{0\}\). Equivalently, the image of restriction from any strictly later level is already \(\{0\}\), so the persistent image is \(\{0\}\). The seed one never extends. This is compatible with KS58, which proves nonemptiness and characterizes the persistent image, not surjectivity of every evaluation.

\subsection*{2. A residual chart can change without changing the defect}
Let \(d:\R^2\to\R\) be \(d(u,r)=r\). For \(a\in\R\), let \(A_a(r)=(ar,r)\). Compute the exact repair \(P_a\), compare it with \(P_b\), and decide when they commute.

\textbf{Solution.} Since \(dA_a=I\),
\[
 P_a(u,r)=(u-ar,0),\qquad
 P_b(u,r)-P_a(u,r)=((a-b)r,0).
\]
Both fix the horizontal exact axis. Thus \(P_aP_b=P_b\) and \(P_bP_a=P_a\). They commute exactly when they are equal, which occurs exactly when \(a=b\), as can be checked at input \((0,1)\). Every repair has zero residual. The different answers reflect the different right inverses, as in KS65.

\subsection*{3. Pairwise invertibility does not supply a triple law}
Let \(t\) be a nonidentity transposition in \(S_3\). Suppose three proposed overlap comparisons require frames \(b_0,b_1,b_2\) satisfying
\[
 b_1=tb_0,\qquad b_2=tb_1,\qquad b_2=tb_0.
\]
Every comparison is individually invertible. Can the three frames exist?

\textbf{Solution.} The first two equations give \(b_2=t^2b_0=b_0\). The third then gives \(b_0=tb_0\), and right cancellation gives \(t=1\), contrary to the assumption. Invertibility of every proposed overlap map does not prove the composition law needed for descent. KS34 applies to data satisfying that law, not to arbitrary lists of pairwise isomorphisms.

\subsection*{4. Zero noisy output is not a certificate of reference exactness}
Take \(D(a,b)=a\), \(\widetilde D(a,b)=a+\tfrac1{10}b\), and input \(x=(-10,100)\). The measured operator returns zero exactly. Find the distance to the reference kernel and evaluate the reference uncertainty bound.

\textbf{Solution.} The reference kernel is the vertical axis, so the distance is ten. Its residual gain is one, and the operator error norm is \(1/10\). With \(y=0\) and \(\eta=0\), KS69 gives
\[
 \dist(x,\ker D)\le\frac1{10}\sqrt{10^2+100^2}=\sqrt{101}.
\]
The true distance ten is below this bound. A formula omitting the operator-error term would instead claim a zero distance, which is false. The input is exact for the measured model and not for the reference one.

\subsection*{5. A prefix computation can hide a wraparound defect}
On a four-cycle, take proposed differences \((2,-1,3,-3)\) and anchor zero. Integrate the first three edges. Which equation fails, and what changes when only the final edge is allowed to move?

\textbf{Solution.} The first three edges give values \((0,2,1,4)\). Their required last difference is \(x_0-x_3=-4\), not \(-3\). The total prescribed sum is one, so the original loop has no realization. Keeping the first three edges forces the unique correction \((2,-1,3,-4)\), whose sum is zero. This is a protected repair of the edge data. It is not a realization of the original data without modification.

\subsection*{6. Residual decay can hide unprotected drift}
In maximum-norm $\R^3$, use $p(s,u,r)=s$ and $d(s,u,r)=r$. Compare the paths $(1,0,(n+1)^{-1})$ and $(1,H_n,(n+1)^{-1})$. Which hypotheses distinguish them?

\textbf{Solution.} Both preserve $p=1$ and have the same residual ratios $(n+1)/(n+2)$. The first moves by $1/((n+1)(n+2))$ and converges to $(1,0,0)$. The second moves by $1/(n+1)$ and has an unbounded second coordinate. In the variable-rate theorem, the first has $C_n=1/(n+2)$ and a summable weighted budget. The second has $C_n=1$, making that budget harmonic and divergent. The difference is movement, not residual decay or semantic preservation.

\subsection*{7. How a finite observation changes the apparent rate}
Let $\eta_n=1/(n+1)$. Find its profile after sampling at $\phi(k)=2^k-1$. Does the new profile prove exponential convergence in the original index?

\textbf{Solution.} The sampled value is $\eta_{\phi(k)}=2^{-k}$. Recovering the original index gives $k=\log_2(n+1)$ on the sampled stages, so the value is still $1/(n+1)$. The schedule preserves vanishing but changes the meaning of the horizontal axis. No faster original process has been proved. KS75 is a convergence equivalence under a monotone cofinal schedule, not a rate-invariance theorem.

\subsection*{8. What does a zero floor fail to certify?}
For the diagonal model of Section~\ref{sec:escape}, can an exact protected answer exist in $\ell^2\oplus\R$? Does a coherent family of exact finite prefixes exist? Explain why the answers differ without contradicting finite-fiber existence.

\textbf{Solution.} Exactness with protected value one forces every sequence coordinate to equal minus one. This is not square summable, so there is no protected Hilbert answer. Each finite prefix nevertheless has the unique solution $(-1,\ldots,-1;1)$, and these solutions are coherent under truncation. Their inverse limit contains the all-minus-one sequence in the full product. The finite-fiber theorem supplies that coherent point. It does not prove that the point is in the image of the square-summable semantic space. The correction image is dense but does not contain the raw residual, so its floor is zero and its algebraic class is nonzero.

\subsection*{9. A phase witness without enumerating every gauge}
For $c(x,y)=x_2y_1$ on $\mathbf F_2^2$, show that no scalar coboundary equals $c$. Then compute the square of the lifted label $(0,e_1+e_2)$ in the eight-element extension.

\textbf{Solution.} Scalar coboundaries $a(x)+a(y)-a(x+y)$ are symmetric, but $c(e_1,e_2)=0$ and $c(e_2,e_1)=1$. Thus no scalar coboundary equals $c$. For $x=e_1+e_2$, $c(x,x)=1$, so $(0,x)^2=(1,0)$ and its fourth power is the identity. The corresponding matrix is $XZ$, whose square is $-I$. The additional phase bit records precisely the sign lost by projecting back to the four-label group.

\subsection*{A checklist for applying the theory}
Before selecting a theorem, state the raw and exact objects, the coefficient structure, the maps that relate views, and the data that must remain fixed. Then ask which conclusion is needed: existence, a description of all answers, an identification-preserving comparison, a stable distance estimate, or an executable finite certificate. Finally record the hypotheses that make that conclusion available.

The examples above fail only when one of those distinctions is suppressed. The missing information is not supplied by giving the same construction a more general name.

\subsection*{Problem 10. The finite energy identity}
In $\R^2$, first prescribe $a_1+a_2=1$, then also prescribe $a_1=3/4$. Find the two minimum norm corrections and check the exact nested-affine identity.

\textbf{Solution.} The first correction is $(1/2,1/2)$, with energy $1/2$. The second is $(3/4,1/4)$, with energy $5/8$. Their difference has squared norm $1/8=5/8-1/2$. The identity uses the affine nature of both constraint sets. For the nested closed convex sets $[1,\infty)$ and $[2,\infty)$ in $\R$, the minimum points instead have squared difference one and energy difference three.

\subsection*{Problem 11. A basis is not automatically an isometry}
Let $p(a)=a_1-2a_2$, $d(a)=a_1$, $r=1$, and $\alpha=1$. Compute the protected regularized correction using the basis $(2,1)$ of the kernel. What changes if the Gram matrix is omitted?

\textbf{Solution.} The coordinate $c$ gives $a=(2c,c)$, with $S=2$ and $G=5$. The normal equation is $9c=2$, so $a=(4/9,2/9)$. Its objective is $5/9$. Replacing $G$ by one gives $c=2/5$ and $a=(4/5,2/5)$, whose actual ambient objective is $21/25>5/9$. That replacement solves a different metric problem.

\subsection*{Problem 12. A robust negative certificate}
A nominal scalar system has $T=1$, $r=1$. The operator error is at most $1/10$ and the data error at most $1/20$. Can a correction of norm at most $1/2$ achieve true residual at most $3/10$? What does the same test say at budget two?

\textbf{Solution.} At budget $1/2$, the nominal optimum is $1/2$. Subtracting $1/20+(1/10)(1/2)$ leaves $2/5>3/10$, so the proposed accuracy is excluded for every admissible true model. At budget two, the lower estimate is negative and clips to zero. It excludes nothing. The disappearance of this certificate is not proof of exact feasibility.

\clearpage
\chap{A hypothesis atlas}\label{app:hypotheses}
A longer monograph is useful only if the reader can locate the theorem that answers the actual question. This atlas groups the main obligations by the desired conclusion. The references lead to the statements and proofs, not just to examples that resemble the intended application.

\subsection{Existence and identification}
\begin{longtable}{@{}p{.24\linewidth}p{.48\linewidth}p{.21\linewidth}@{}}
\toprule Desired conclusion & Data and hypotheses to establish & Main entry\\\midrule\endhead
A coherent family & A view system with its actual restrictions and all compatibility equations. Nonemptiness is separate. & KS01--KS03\\[4pt]
A prescribed extension & The declared partial section and an actual vanishing relative obstruction, or matching data on a transition cover. & KS04, KS09, KS14\\[4pt]
A unique extension & Existence and vanishing of the invisible coherent ambiguity, not existence alone. & KS10, KS32\\[4pt]
An infinite point from finite views & Nonempty finite fibers, or another proved lifting/compactness hypothesis. A prescribed value must lie in the persistent image. & KS54, KS58\\[4pt]
Descent of objects and witnesses & The actual groupoid data, coherent overlap isomorphisms, and a composition-covering family on the specified site. & KS34--KS36\\[4pt]
Invariance of a boundary problem & An identified boundary and the specified comparison equivalences. Literal freezing of every representative is a different construction. & KS38--KS39, KS68\\[4pt]
The same limit from two paths & Convergence of both paths in a metric space and a pairwise distance budget tending to zero. & KS74\\
\bottomrule
\end{longtable}

\subsection{Approximation and numerical meaning}
\begin{longtable}{@{}p{.24\linewidth}p{.48\linewidth}p{.21\linewidth}@{}}
\toprule Desired conclusion & Data and hypotheses to establish & Main entry\\\midrule\endhead
An exact limit of approximations & Completeness, a nonnegative lower semicontinuous residual, vanishing residuals, and summable candidate movement. & KS71--KS72\\[4pt]
A protected exact limit & In addition, continuity of the protected map and preservation along the candidate path. & KS73\\[4pt]
Exact cancellation from zero floor & Closedness of the actual admissible residual image. Without it only closure membership follows. & KS76--KS77\\[4pt]
A stable residual-to-distance bound & The required closed range or bounded inverse in the declared normed model. In a protected problem, use the restriction to $\ker p$. & KS43, KS56\\[4pt]
A controlled moving exact space & Two residual certificates, or a proved split-surjective regime with the strict smallness condition. Norm closeness alone does not preserve rank. & KS48--KS53\\[4pt]
An independent object distance & An arrow category and an extended nonnegative cost with zero identity cost and composition subadditivity. Separation is another question. & KS78--KS79\\
\bottomrule
\end{longtable}

\subsection{Bounded budgets and regularization: written extensions}
The following entries refer to full printed proofs, not new Lean declarations.
\begin{longtable}{@{}p{.24\linewidth}p{.48\linewidth}p{.21\linewidth}@{}}
\toprule Desired conclusion & Data and hypotheses to establish & Main entry\\\midrule\endhead
An exact answer from bounded approximations & A fixed bounded linear equation between real Hilbert spaces, residual convergence, and a uniform correction bound. & WR01\\[4pt]
Strong convergence of finite repairs & Nested nonempty closed affine constraints, minimum norm answers, and a uniform energy budget. & WR02\\[4pt]
A sharp finite-budget exclusion & The correction ball, the restricted adjoint, and a unit dual witness. & WR03--WR04\\[4pt]
Robustness to model uncertainty & A fixed correction space and certified operator and data error bounds. & WR05\\[4pt]
A unique protected regularized answer & Positive regularization parameter and the actual Hilbert metric on the protected kernel. & WR06--WR07\\[4pt]
A convergent regularization limit & Actual attainment of the projected datum, not merely its membership in the closure. & WR08\\[4pt]
A noisy convergent answer & The same attained datum, vanishing parameter, and noise divided by the square root of the parameter tending to zero. & WR09, WR12\\[4pt]
A leading approximation-cost constant & The explicit diagonal model and a proof controlling its infinite Riemann-sum tails. & WR10--WR11\\\bottomrule
\end{longtable}

\subsection{Three questions before using a finite check}
First determine whether the tested object is the actual object in the theorem. A scalar simulation can illustrate an operator-valued statement without checking every dependent fiber. Next determine whether the test checks a universal finite domain or a selected finite range. Finally identify any limit that remains outside the computation.

The new phase extension is exhaustively checked on its eight elements. The diagonal example is different: rational finite sums verify its displayed enclosures, while the proof of nonexistence uses the infinite square-summability condition. Calling both results simply ``tested'' would conceal the most important distinction.

The companion evidence guide, \path{data/deepening_models.json}, and \path{data/written_extensions.json} retain these scopes. The general source declarations are checked in the accepted formal baseline. Printed examples and new conceptual connections do not acquire additional Lean coverage merely because they occur next to a source identifier.

\chap{Notation and source concordance}\label{app:concordance}
\subsection{Notation}
\begin{longtable}{@{}p{.25\linewidth}p{.7\linewidth}@{}}
\toprule
Symbol & Meaning\\\midrule\endhead
\(\V,F_v,f^*\) & A declared view category, its fibers, and contravariant restrictions.\\
\(\Gamma(F)\) & The strict set of coherent global families.\\
\(\Gamma_U(F),\Ext_U(s)\) & Partial sections on a full patch and their global extension fiber.\\
\(\delta,\partial,d,D\) & Respectively the view differential, tower differential, a model residual, and a Hilbert defect operator. Each target is declared where used.\\
\(p,p_U,p_L\) & A protected observable, restriction to prescribed views, and a splitting dependent innovation observable.\\
\(\mathfrak o_U,H_U\) & The existence obstruction and the group of invisible coherent changes. Additive and nonabelian conventions are stated separately.\\
\(C_j,J_j,Q_N,S_N\) & Coarsener, chosen lift, restored coarse remainder, and retained finite reconstruction.\\
\(c_f,a_v\) & Coherent nonabelian transition labels and gauge values.\\
\(\mathfrak K_G,J_{\rm comp}\) & The cocycle patch groupoid and its composition-support topology.\\
\(\mathcal E_b(R)\) & Realizations equipped with an actual identification of their boundary with \(b\).\\
\(T_i,K_{ij},R_p\) & Transport to the root, the corresponding root loop operator, and reconstruction from the root.\\
\(\mathcal H,q,P\) & Root residual, an orthogonal kernel projector, and the explicitly defined model repair. A graph repair need not be orthogonal in the global maximum norm.\\
\(\pi_{W,U},H_U\) & The map forgetting part of the prescription and the corresponding invisible coherent ambiguity.\\
\(\phi,\operatorname{EventualHom}\) & A monotone source schedule and transport classes modulo eventual equality.\\
\(\chi_A,M=SA-BC\) & Split exact/residual coordinates and the mismatch between chosen repair lifts.\\
\(\omega,Q_n,C_n,b_n\) & A residual, its prefix-product rate, movement coefficient, and summable step budget. Here $Q_n$ is not the dyadic remainder operator $Q_N$.\\
\(\nu_B,d_W,d_W^{\mathrm{sym}}\) & Residual floor, extended directed transport distance, and its symmetric maximum.\\
\(c(x,y),\kappa(x,y)\) & The two-bit pair phase and its antisymmetric scalar-gauge invariant.\\
\(H,T,r\) & In the budget chapters, $H=\ker p$ with its inherited Hilbert metric, $T=d|_H$, and $r=df$.\\
\(\beta_{T,r}(M),a_\alpha\) & Minimum residual at correction budget $M$, and the quadratic regularized correction.\\
\(C,b,q\) & Closure of the correction image, its projection of the residual, and the irreducible residual floor.\\
\(\simeq,\simeq_L\) & A specified equivalence, and a continuous linear equivalence. Neither notation alone asserts an isometry.\\
\bottomrule
\end{longtable}

\subsection{How to read the correspondence}
The identifiers below are stable editorial identifiers, not counts of new discoveries. A row can combine several source declarations into one exposition. The named declaration, its line range, source hash, and correspondence note are recorded in \code{data/claim_ledger.json} and in the companion Lean guide. Existence of a matching source declaration is not a mechanical proof that every word of an edited statement matches its type. The displayed hypotheses and the printed proof remain part of the review.

\textbf{F} denotes the recovered foundation tree at version 0.20.0. \textbf{M} denotes the focused moving-space tree at version 0.32.0. \textbf{S} denotes the separately recovered Hilbert holonomy source. \textbf{I} denotes a source in the canonical integrated tree, whether inherited or added during the version 1.2.0 formalization. F, M, and S retain provenance, while byte-identical copies of their referenced sources were compiled together in the accepted version 1.2.0 replay of \path{formal/KomorebiSpace/}. All \EditionClaims\ KS catalogue items have source correspondences in that checked closure. The twelve WR extensions are separate written results, listed in Appendix~\ref{app:written}, and have no new Lean counterparts in this edition. This does not assert that every illustrative computation or explanatory sentence is independently formalized.

\input{claim_table}

\par\noindent\begin{minipage}{\linewidth}\subsection{Verification is a separate record}
The foundation archive reports 122 production proof modules, 1,299 named theorem declarations, and 696 indexed constructions. Its recorded strategy includes one complete cumulative source build, followed by a fresh archive replay that reuses a checked 117-module baseline and compiles the five additions and root again. It must not be described as two full clean cumulative builds.

The focused moving-space archive reports 11 production modules and a root, 151 named theorem declarations, and 47 indexed constructions. Its historical replay record reports two complete clean builds of that focused source closure. The archive SHA-256 and the retained source hashes are checked by this publication's integrity script. The focused tree does not replace the whole historical Komorebi Space repository.

The pinned environment in the records is Lean 4.34.0-rc2 with mathlib commit\par\smallskip\noindent\path{aeb72f1e9981b6d35b426b16dca7759be3e7ba06}.\par\smallskip The permitted standard axioms are \code{propext}, \code{Classical.choice}, and \code{Quot.sound}. The compiler, kernel, and cached dependencies remain trusted. For version 1.2.0 the portable runtime was restored and its version and dependency pins checked, but the compiler and mathlib were not rebuilt independently. This expanded publication retains that verified baseline without executing a new compiler replay.

The retained version 1.2.0 formal record covers \EditionFormalModules\ production modules and one root. Clean source replays and an imported axiom audit are separate from archive integrity, exact finite computations, and PDF preflight. The five modules added in that formalization contain \EditionNewTheorems\ explicit named theorems and lemmas, including supporting results. Counts are not claims of mathematical novelty. The complete source hashes, command logs, and baseline scope are in \path{audit/formal/release.json}.

\end{minipage}

\subsection*{Editorial source policy}
Earlier source documents contain historical scope statements and counts. Later checked developments can supersede an earlier statement that a topic is a target. In particular the old outline's stack target is not the status of the later standard stack theorem. This monograph uses the actual relevant declarations rather than treating the earliest outline as a current completion ledger.



\subsection*{Retained source connections and the written extension}
The earlier expansion developed KS59--KS70. KS71--KS80 add exactification, floors, transport costs, and the phase obstruction. Their declarations are listed above. Separately scoped worked models appear in \path{data/deepening_models.json}. WR01--WR12 are indexed separately in \path{data/written_extensions.json}. They have complete printed proofs and exact finite checks, but no new compiled declarations. No new compiler execution is claimed.

\clearpage
\chap{The bounded-budget extension: proofs and evidence}\label{app:written}
WR01--WR12 are a deliberately separate catalogue. Their complete proofs are printed in Sections~\ref{sec:bounded-realization}--\ref{sec:diagonal-regularization}. None is counted among the eighty source-linked KS items or the retained Lean declarations. The new exact checker audits finite identities, rational programs, and enclosures. It does not supply machine proofs of weak compactness, infinite convergence, or the asymptotic integrals.
\begin{longtable}{@{}p{.10\linewidth}p{.38\linewidth}p{.43\linewidth}@{}}
\toprule ID & Result & Hypotheses that matter\\\midrule\endhead
\hyperref[WR01]{WR01} & Bounded approximate feasibility & Real Hilbert spaces, bounded linear T, one fixed residual, uniform norm bound on approximate corrections.\\[5pt]
\hyperref[WR02]{WR02} & Strong convergence of minimum energy finite repairs & Nested nonempty closed affine sets in one Hilbert space. Minimum norm corrections. Uniform budget is equivalent to a nonempty intersection.\\[5pt]
\hyperref[WR03]{WR03} & Attained budget residual profile & Bounded linear T between real Hilbert spaces, fixed r and finite nonnegative budget.\\[5pt]
\hyperref[WR04]{WR04} & Sharp dual budget certificate & Same fixed Hilbert model. Use the actual restricted adjoint. Supports zero budget and zero residual.\\[5pt]
\hyperref[WR05]{WR05} & Robust budget profile & Operators on the same fixed correction space. Norm errors eta and delta; no uncertainty in the protected kernel.\\[5pt]
\hyperref[WR06]{WR06} & Protected regularized minimizer and energy identity & Positive alpha, bounded linear T, actual Hilbert metric on ker p. No closed-range or compact-operator assumption.\\[5pt]
\hyperref[WR07]{WR07} & Parameter monotonicity and budget optimality & Same operator and residual at two positive regularization parameters.\\[5pt]
\hyperref[WR08]{WR08} & The regularization limit alternative & b is the projection onto closure of range. Strong convergence requires b in actual range; otherwise correction norms diverge.\\[5pt]
\hyperref[WR09]{WR09} & Noise sensitivity and parameter choice & Fixed T and positive alpha. For convergence, attainable projected exact datum, alpha tends to zero, and delta/sqrt(alpha) tends to zero.\\[5pt]
\hyperref[WR10]{WR10} & Diagonal correction and infinite-tail enclosures & $T=\operatorname{diag}(1/n)$, $r=(1/n)$, integer $m\ge1$, $\alpha=m^{-2}$, cutoff $K\ge m$.\\[5pt]
\hyperref[WR11]{WR11} & Sharp asymptotic budget cost & The same fixed infinite diagonal problem. Riemann sums justified with uniform tail bounds. Not a universal efficiency factor.\\[5pt]
\hyperref[WR12]{WR12} & Explicit noisy attainable datum & The attainable datum $(1/n^2)$, deterministic noise $0<\delta\le1$, and $m=\lceil\delta^{-2/3}\rceil$.\\[5pt]
\bottomrule\end{longtable}
\subsection{Independent finite checks}
The regularization suite runs \EditionRegularizationChecks\ exact assertions. Its direct two-dimensional oracle uses the adjugate formula, independently of the program's elimination routine. The protected problems are also compared with full ambient Karush--Kuhn--Tucker systems, avoiding the program's nullspace coordinate route. The executable program retains a nonorthogonal kernel basis's Gram matrix. Input validation rejects zero regularization, ragged dimensions, floating-point rationals, and invalid tail cutoffs.

The normal-equation residual, preservation values, correction norm, and fit residual are returned explicitly. The diagonal demonstration prints a finite sum together with a proved enclosure of its omitted infinite tail. Those enclosure statements are WR10's printed inequalities. They are not inferred from the finite windows used by the checker.

\subsection{Formal boundary and restart}
The canonical formal tree is unchanged. Formalizing these extensions requires new source modules, explicit comparison to the existing protected-repair APIs, and an actual clean compiler and axiom replay. The relevant starting points are the inherited Hilbert projection and protected correction modules. A proof outline or an import-only file must not be counted as formal coverage. The machine-readable statement ledger is \path{data/written_extensions.json}.


\clearpage
\chap{Background and source documents}\label{app:references}
The references identify the standard background and frozen project sources. Printed arguments are supplied for the principal constructions. The dyadic reconstruction argument is now written out, including its starting-scale induction. The five additions KS54 through KS58 now have corresponding declarations in the integrated formal project. Their source notes distinguish the formal proof route from the printed exposition. Live mathlib documentation is a navigation aid, not a replacement for the pinned source revision.

\begin{thebibliography}{99}
\bibitem{stacks-limits} The Stacks Project Authors. \emph{Limits and colimits over preordered sets}. Section 4.21, Tag 002Z, including the finite-set inverse-limit theorem in Lemma 4.21.7. \url{https://stacks.math.columbia.edu/tag/002Z}. Accessed 10 October 2026.
\bibitem{stacks-stack} The Stacks Project Authors. \emph{Stacks}. Section 8.4, Tag 0268. In particular the effective descent and equivalence-of-descent-categories formulations. \url{https://stacks.math.columbia.edu/tag/0268}. Accessed 10 October 2026.
\bibitem{mathlib-stack} The mathlib Community. \emph{CategoryTheory.Sites.Descent.IsStack}. Standard pseudofunctor stack interface. \url{https://leanprover-community.github.io/mathlib4_docs/Mathlib/CategoryTheory/Sites/Descent/IsStack.html}.
\bibitem{mathlib-projection} The mathlib Community. \emph{Analysis.InnerProductSpace.Projection.Basic}. Orthogonal projection onto complete subspaces. \url{https://leanprover-community.github.io/mathlib4_docs/Mathlib/Analysis/InnerProductSpace/Projection/Basic.html}.
\bibitem{mathlib-banach} The mathlib Community. \emph{Analysis.Normed.Operator.Banach}. Banach open mapping and continuous inverse constructions. \url{https://leanprover-community.github.io/mathlib4_docs/Mathlib/Analysis/Normed/Operator/Banach.html}.
\bibitem{foundation} Komorebi P. \emph{Komorebi Space Theory}, version 0.20.0. Recovered source archive and replay evidence, 28 September 2026. Included unchanged under \code{evidence/archives/}. Its source documents include \emph{Ambient Sections and Extension Obstructions}, \emph{Split Tower Coordinates}, \emph{Infinite Reconstruction}, \emph{Nonabelian Coherence}, \emph{Effective Groupoid Descent}, \emph{A Stack on the Composition Site}, and \emph{Nonisometric Transport and Amplified Reconstruction}.
\bibitem{moving} Komorebi P. \emph{Komorebi Space: Controlled Motion of Exact Realization Spaces}, focused version 0.32.0. Historical source archive and replay evidence, 30 September 2026. Included unchanged under \code{evidence/archives/}.
\bibitem{hilbert-extra} Komorebi P. \emph{HilbertHolonomyRepair.lean}. Separately recovered source implementing closed-range amplified holonomy reconstruction. Included under \code{evidence/selected/}. Its original separate provenance is retained. A byte-identical canonical copy was subsequently compiled with its dependencies in the accepted version 1.2.0 integrated replay.
\bibitem{clason-inverse} Christian Clason. \emph{Regularization of Inverse Problems}. Graduate lecture notes, arXiv:2001.00617v2, 2021. Sections 2, 4, and 6 provide standard background on weak convergence, parameter choice, and Tikhonov regularization. \url{https://arxiv.org/abs/2001.00617v2}. Consulted 12 October 2026.
\end{thebibliography}


\end{document}
